\documentclass[11pt,a4paper]{article}
\usepackage[utf8]{inputenc}
\usepackage[vietnamese]{babel}
\usepackage{amsmath,amssymb}
\usepackage{geometry}
\geometry{a4paper, total={180mm,260mm}, left=15mm, right=15mm, top=18mm, bottom=18mm}

\usepackage{tikz}
\usetikzlibrary{shapes.geometric, arrows.meta, positioning, calc}
\tikzset{arr/.style={-{Stealth[scale=0.8]}}}
\usepackage{tcolorbox}
\tcbuselibrary{skins,breakable}
\usepackage{fontawesome}
\usepackage{fancyhdr}
\usepackage{enumitem}

\pagestyle{fancy}
\fancyhf{}
\renewcommand{\headrulewidth}{0.5pt}
\renewcommand{\footrulewidth}{0.5pt}
\fancyhead[L]{\small\bfseries\color{violet!85!black}\faBook\ SỔ TAY KIẾN THỨC TOÁN 2}
\fancyhead[R]{\small\bfseries\color{violet!85!black}Lớp Toán Thầy Lĩnh -- loptoanthaylinh.vn}
\fancyfoot[C]{\small\color{gray}-- Sổ tay kiến thức: Bài 03. Các thành phần của phép cộng, phép trừ --}

\setlength{\parindent}{0pt}
\setlength{\parskip}{0pt}

\begin{document}

% --- BANNER BÀI HỌC ---
\begin{tcolorbox}[
  enhanced,
  colback=violet!8!white,
  colframe=violet!70!black,
  arc=3mm, boxrule=1.2pt,
  left=3mm, right=3mm, top=2mm, bottom=2mm,
  halign=center
]
  {\LARGE\bfseries\color{violet!90!black}\faBookmark\ BÀI 03: CÁC THÀNH PHẦN CỦA PHÉP CỘNG, PHÉP TRỪ}\\[1mm]
  {\itshape\footnotesize\color{black!75} Tên gọi các thành phần và kết quả của phép cộng, phép trừ; ý nghĩa danh xưng trên 1 trang A4}
\end{tcolorbox}

\vspace{-2mm}

\noindent
\begin{minipage}[t]{0.49\textwidth}\vspace{0pt}
\begin{tcolorbox}[
  enhanced,
  colback=violet!4!white,
  colframe=violet!75!black,
  arc=2mm, boxrule=0.8pt,
  left=2mm, right=2mm, top=1.8mm, bottom=1.8mm,
  title={\footnotesize\bfseries\color{white}\faKey\ 1. Thành phần phép cộng}
]
{\footnotesize
\begin{itemize}\setlength{\itemsep}{1pt}\setlength{\leftskip}{-2.5mm}
  \item \textbf{Cấu trúc tổng quát:}
  $$\mathbf{a + b = c}$$
  \begin{itemize}\setlength{\itemsep}{0.5pt}\setlength{\leftskip}{-2.5mm}
    \item $a$: \textbf{Số hạng thứ nhất}.
    \item $b$: \textbf{Số hạng thứ hai}.
    \item $c$: \textbf{Tổng}.
  \end{itemize}
  \item \textbf{Lưu ý đặc biệt:}\\
  Biểu thức $a + b$ cũng được gọi là một \textbf{"tổng"}.\\
  \textit{Ví dụ:} Tổng của $24$ và $15$ là $24 + 15$ hay $39$.
  \item Khi đổi chỗ các số hạng thì tổng không thay đổi: $a + b = b + a$.
\end{itemize}
}
\end{tcolorbox}
\end{minipage}
\hfill
\begin{minipage}[t]{0.49\textwidth}\vspace{0pt}
\begin{tcolorbox}[
  enhanced,
  colback=teal!4!white,
  colframe=teal!75!black,
  arc=2mm, boxrule=0.8pt,
  left=2mm, right=2mm, top=1.8mm, bottom=1.8mm,
  title={\footnotesize\bfseries\color{white}\faSliders\ 2. Thành phần phép trừ}
]
{\footnotesize
\begin{itemize}\setlength{\itemsep}{1pt}\setlength{\leftskip}{-2.5mm}
  \item \textbf{Cấu trúc tổng quát:}
  $$\mathbf{a - b = c}$$
  \begin{itemize}\setlength{\itemsep}{0.5pt}\setlength{\leftskip}{-2.5mm}
    \item $a$: \textbf{Số bị trừ} (đứng trước dấu trừ).
    \item $b$: \textbf{Số trừ} (đứng sau dấu trừ).
    \item $c$: \textbf{Hiệu} (kết quả phép trừ).
  \end{itemize}
  \item \textbf{Lưu ý đặc biệt:}\\
  Biểu thức $a - b$ cũng được gọi là một \textbf{"hiệu"}.\\
  \textit{Ví dụ:} Hiệu của $58$ và $23$ là $58 - 23$ hay $35$.
  \item Trong số tự nhiên, \textbf{Số bị trừ luôn $\ge$ Số trừ}.
\end{itemize}
}
\end{tcolorbox}
\end{minipage}

\vspace{-1.5mm}
\begin{tcolorbox}[
  enhanced,
  colback=orange!4!white,
  colframe=orange!80!black,
  arc=2mm, boxrule=0.8pt,
  left=2mm, right=2mm, top=1.5mm, bottom=1.5mm,
  title={\footnotesize\bfseries\color{white}\faSitemap\ 3. Bảng đối chiếu thành phần phép cộng và phép trừ}
]
\centering
\begin{tikzpicture}[
  box/.style={rectangle, rounded corners=2pt, draw=#1!80!black, fill=#1!10, text width=2.6cm, align=center, font=\scriptsize\bfseries, inner sep=3pt},
  arr/.style={-{Stealth[scale=0.8]}, thick, black!60}
]
  % Phép cộng
  \node[box=violet] (c1) {SỐ HẠNG\\($32$)};
  \node[right=0.2cm of c1, font=\normalsize\bfseries] (p1) {$+$};
  \node[box=violet, right=0.2cm of p1] (c2) {SỐ HẠNG\\($14$)};
  \node[right=0.2cm of c2, font=\normalsize\bfseries] (eq1) {$=$};
  \node[box=green, right=0.2cm of eq1] (c3) {TỔNG\\($46$)};

  % Phép trừ
  \node[box=teal, below=0.6cm of c1] (t1) {SỐ BỊ TRỪ\\($58$)};
  \node[right=0.2cm of t1, font=\normalsize\bfseries] (m1) {$-$};
  \node[box=teal, right=0.2cm of m1] (t2) {SỐ TRỪ\\($20$)};
  \node[right=0.2cm of t2, font=\normalsize\bfseries] (eq2) {$=$};
  \node[box=orange, right=0.2cm of eq2] (t3) {HIỆU\\($38$)};
\end{tikzpicture}
\end{tcolorbox}

\vspace{-1.5mm}
\noindent
\begin{minipage}[t]{0.49\textwidth}\vspace{0pt}
\begin{tcolorbox}[
  enhanced,
  colback=green!4!white,
  colframe=green!60!black,
  arc=2mm, boxrule=0.8pt,
  left=2mm, right=2mm, top=1.8mm, bottom=1.8mm,
  title={\footnotesize\bfseries\color{white}\faCheckCircle\ 4. Ví dụ thực hành chuẩn mực}
]
{\footnotesize
\textbf{1. Gọi tên thành phần:}\\
- Trong phép tính $45 + 32 = 77$:\\
  $45$ và $32$ là các số hạng; $77$ là tổng.\\
- Trong phép tính $89 - 54 = 35$:\\
  $89$ là số bị trừ; $54$ là số trừ; $35$ là hiệu.
\textbf{2. Đặt tính theo yêu cầu:}\\
Biết số bị trừ là $67$, số trừ là $23$. Tính hiệu?\\
\textit{Lời giải:} Hiệu là: $67 - 23 = \mathbf{44}$.
}
\end{tcolorbox}
\end{minipage}
\hfill
\begin{minipage}[t]{0.49\textwidth}\vspace{0pt}
\begin{tcolorbox}[
  enhanced,
  colback=red!4!white,
  colframe=red!75!black,
  arc=2mm, boxrule=0.8pt,
  left=2mm, right=2mm, top=1.8mm, bottom=1.8mm,
  title={\footnotesize\bfseries\color{white}\faExclamationTriangle\ 5. Mẹo tránh bẫy trắc nghiệm}
]
{\footnotesize
\begin{itemize}\setlength{\itemsep}{1pt}\setlength{\leftskip}{-2.5mm}
  \item \textbf{Bẫy từ khóa "Tính tổng" vs "Tính hiệu":}\\
  Đề bài: "Tính tổng của 30 và 20" $\implies$ Làm tính \textbf{cộng} ($30 + 20 = 50$). Không được làm trừ!
  \item \textbf{Bẫy phân biệt Số bị trừ và Số trừ:}\\
  Số đứng \textbf{trước} dấu trừ là \textbf{Số bị trừ} (bị bớt đi), số đứng \textbf{sau} dấu trừ là \textbf{Số trừ}.
  \item \textbf{Bẫy biểu thức là một tổng/hiệu:}\\
  $25 + 14$ cũng chính là một tổng!
\end{itemize}
}
\end{tcolorbox}
\end{minipage}


\end{document}
